\chapter{Непрерывность преобразования Фурье абсолютно интегрируемой функции. Преобразование Фурье производной и производная преобразования Фурье.}

\section{Непрерывность преобразования Фурье абсолютно интегрируемой функции}

\begin{defn}
Определённая на $\bbR$ функция называется \textit{локально интегрируемой}\rindex{функция!локально интегрируемая}, если она абсолютно интегрируема на любом конечном интервале.
\end{defn}

\begin{defn}
Для любой локально интегрируемой функции $\phi(x), x \in \bbR$ предел
$$
\lim_{l \to +\infty}\int_{-l}^{l} \phi(x)\,dx
$$
называют \textit{интегралом от $-\infty$ до $+\infty$ в смысле главного значения (или в смысле Коши)}\rindex{интеграл!в смысле главного значения} и обозначают
$$
\v.p.\int_{-\infty}^{+\infty} \phi(x)\,dx.
$$
\end{defn}

Аналогом ряда Фурье в данном вопросе будем называть интеграл
$$
\int_{0}^{+\infty} (a(y)\cos(yx) + b(y)\sin(yx))\,dy,\, \text{где}
$$
$$
a(y) = \lim_{l \to +\infty}\frac{1}{\pi}\int_{-l}^{l} f(t)\cos(yt)\,dt;\qquad 
b(y) = \lim_{l \to +\infty}\frac{1}{\pi}\int_{-l}^{l} f(t)\sin(yt)\,dt
$$
который называется \textit{интегралом Фурье}\rindex{интеграл!Фурье}.

Заметим, что не для всякой функции $f$, которая абсолютно интегрируема на любом конечном интервале, определённые выше пределы существуют. Если же функция $f$ абсолютно интегрируема на $\bbR$, то эти пределы заведомо существуют и
\begin{equation} \label{ch19eq1}
a(f;y) = \frac{1}{\pi}\int_{-\infty}^{+\infty} f(t)\cos(yt)\,dt;
\qquad
b(f;y) = \frac{1}{\pi}\int_{-\infty}^{+\infty} f(t)\sin(yt)\,dt.
\end{equation}

\begin{thm}[лемма Римана об осцилляции\rindex{лемма!Римана об осцилляции}]
\label{th22:1}
Если $f \in L_R(I)$, где $I$~--- произвольный промежуток в $\bbR$, то
$$
\lim_{t \to \infty}\int_{I} f(x)\cos(tx)\,dx = 0, \quad
\lim_{t \to \infty}\int_{I} f(x)\sin(tx)\,dx = 0.
$$
\end{thm}

\begin{proof}
Проведём доказательство для синуса, для косинуса аналогично.

Рассмотрим сначала случай, когда $I = [a;b]$~--- конечный отрезок, а функция $f$ интегрируема по Риману на $[a;b]$. По пункту~2 критерия Дарбу
$$
\forall\, \varepsilon > 0 \quad \exists\, R \text{ (разбиение $[a;b]$)}: \quad
\sum_{i=1}^{N} (M_i - m_i)\Delta x_i < \frac{\varepsilon}{2}.
$$
Тогда
\begin{multline*}
\left| \int_{a}^{b} f(x)\sin(tx)\,dx \right|
= \left| \sum_{i=1}^{N} \int_{x_{i-1}}^{x_i} f(x)\sin(tx)\,dx \right| = \\
= \left| \sum_{i=1}^{N} \int_{x_{i-1}}^{x_i} (f(x) - m_i)\sin(tx)\,dx
       + \sum_{i=1}^{N} m_i \int_{x_{i-1}}^{x_i} \sin(tx)\,dx \right| \le \\
\le \sum_{i=1}^{N} \int_{x_{i-1}}^{x_i} |f(x) - m_i|\,dx
  + \sum_{i=1}^{N} \frac{|m_i|}{|t|} \cdot |\cos(tx_{i-1}) - \cos(tx_i)| \le \\
\le \sum_{i=1}^{N} (M_i - m_i)\Delta x_i + 2N \cdot \frac{M}{|t|},
\quad \text{где } M = \sup_{[a;b]} |f(x)|.
\end{multline*}

Так как числа $M$ и $N$ фиксированы, то
$$
\exists\, t_0: \quad \forall\, t,\ |t| > t_0 \implies \frac{2MN}{|t|} < \frac{\varepsilon}{2}.
$$
Тогда
$$
\forall\, \varepsilon > 0 \quad \exists\, t_0 > 0: \quad
\forall\, t,\ |t| > t_0 \implies
\left| \int_{a}^{b} f(x)\sin(tx)\,dx \right| < \frac{\varepsilon}{2} + \frac{\varepsilon}{2} = \varepsilon,
$$
т.е. $\lim\limits_{t \to \infty} \int_{a}^{b} f(x)\sin(tx)\,dx = 0$.

\smallskip
Пусть теперь $f$ имеет конечное число особенностей на произвольном промежутке $I$, и $\int_I |f(x)|\,dx$ сходится. Не уменьшая общности, можно считать, что особенность одна, причём особенностью является конец промежутка (для определённости~--- правый); иначе $I$ можно разбить на конечное число полуинтервалов, на каждом из которых $f$ имеет единственную особенность в открытом конце, и доказать утверждение для каждого из этих полуинтервалов в отдельности.

Итак, в силу сходимости $\int_a^{\to b}|f(x)|\,dx$,
$$
\forall\, \varepsilon > 0 \quad \exists\, b' \in (a;b):
$$
на $[a;b']$ функция $f$ интегрируема по Риману, а
$$
\int_{b'}^{\to b}|f(x)|\,dx < \frac{\varepsilon}{2}.
$$

Так как для интегрируемой по Риману функции теорема уже доказана,
$$
\exists\, t_0 > 0: \quad \forall\, t,\ |t| > t_0 \implies
\left| \int_{a}^{b'} f(x)\sin(tx)\,dx \right| < \frac{\varepsilon}{2}.
$$

Тогда $\forall\, t,\ |t| > t_0$:
\begin{multline*}
\left| \int_{a}^{\to b} f(x)\sin(tx)\,dx \right|
\le \left| \int_{a}^{b'} f(x)\sin(tx)\,dx \right|
  + \left| \int_{b'}^{\to b} f(x)\sin(tx)\,dx \right| < \\
< \frac{\varepsilon}{2} + \int_{b'}^{\to b}|f(x)|\,dx
< \frac{\varepsilon}{2} + \frac{\varepsilon}{2} = \varepsilon,
\end{multline*}
значит, $\lim\limits_{t \to \infty} \int_a^{\to b} f(x)\sin(tx)\,dx = 0$.

\begin{lemm}[о равномерной непрерывности $a(y)$ и $b(y)$]
\label{lm25:1}
Если $f \in L_R(-\infty, +\infty)$, то функции $a(y)$ и $b(y)$, определённые из~\eqref{ch19eq1}, равномерно непрерывны на $(-\infty; +\infty)$.
\end{lemm}

\begin{proof}
При всех $y_1, y_2 \in \bbR$ имеем
\begin{multline*}
|a(y_2) - a(y_1)|
= \frac{1}{\pi}\left| \int_{-\infty}^{+\infty} f(t)(\cos(ty_2) - \cos(ty_1))\,dt \right| \le \\
\le \frac{1}{\pi} \int_{-\infty}^{+\infty} |f(t)| \cdot 
\left| 2\sin\frac{t(y_1 - y_2)}{2} \sin\frac{t(y_1 + y_2)}{2} \right| dt \le \\
\le \int_{-\infty}^{-A} |f(t)|\,dt + \int_{A}^{+\infty} |f(t)|\,dt
  + \int_{-A}^{A} |tf(t)| \cdot |y_2 - y_1|\,dt,
\end{multline*}
где $A > 0$~--- произвольное число.

Зафиксируем $\varepsilon > 0$. Так как $\int_{-\infty}^{+\infty}|f(t)|\,dt$ сходится, то можно выбрать $A = A(\varepsilon)$ так, что
$$
\int_{A}^{+\infty}|f(t)|\,dt < \frac{\varepsilon}{3}, \quad
\int_{-\infty}^{-A}|f(t)|\,dt < \frac{\varepsilon}{3}.
$$
Пусть $\int_{-A}^{A}|tf(t)|\,dt = M(\varepsilon)$ (если этот интеграл равен $0$, то $M(\varepsilon)$~--- любое положительное число); $\delta(\varepsilon) = \dfrac{\varepsilon}{3M(\varepsilon)}$. Тогда
$$
\forall\, y_1, y_2,\ |y_2 - y_1| < \delta \implies
|a(y_2) - a(y_1)| < \frac{2\varepsilon}{3} + \delta \cdot M(\varepsilon) = \varepsilon.
$$
Значит, функция $a(y)$ равномерно непрерывна на $(-\infty; +\infty)$. Доказательство для функции $b(y)$ аналогично.
\end{proof}

\begin{thm}
\label{th25:4}
Если комплекснозначная функция $f \in L_R(-\infty, +\infty)$, то функции $\widehat{f}(y)$ и $\widetilde{f}(y)$ равномерно непрерывны на $(-\infty, +\infty)$, при этом
$$
\lim_{y \to \infty} \widehat{f}(y) = \lim_{y \to \infty} \widetilde{f}(y) = 0.
$$
\end{thm}

\begin{proof}
Это следует из того, что такими же свойствами обладают и функции $\int_{-\infty}^{+\infty} f(x)\cos(xy)\,dx$ и $\int_{-\infty}^{+\infty} f(x)\sin(xy)\,dx$ (с точностью до числовых множителей действительная и мнимая части $\widehat{f}(y)$ и $\widetilde{f}(y)$)~--- по Лемме~\ref{lm25:1} и лемме Римана.
\end{proof}


\begin{defn}
\label{df25:4}
Пусть $f$~--- комплекснозначная функция действительной переменной, абсолютно интегрируемая на любом конечном отрезке. \textit{Преобразованием (или образом) Фурье функции $f$}\rindex{преобразование Фурье} называется комплекснозначная функция действительной переменной
$$
\widehat{f}(\xi) \equiv F[f](\xi) = \v.p.\frac{1}{\sqrt{2\pi}}\int_{-\infty}^{+\infty} f(x)e^{-i\xi x}\,dx;
$$
\textit{обратным преобразованием (или прообразом) Фурье функции $f$}\rindex{преобразование Фурье!обратное} называется комплекснозначная функция
$$
\widetilde{f}(\xi) \equiv F^{-1}[f](\xi) = \v.p.\frac{1}{\sqrt{2\pi}}\int_{-\infty}^{+\infty} f(x)e^{i\xi x}\,dx
$$
(предполагается, что оба этих интеграла существуют при всех $\xi \in \bbR$).
\end{defn}

\begin{notion}
При всех $\xi \in \bbR$ выполняется равенство $\widetilde{f}(\xi) = \widehat{f}(-\xi)$.

Если $f \in L_R(-\infty, +\infty)$, то все эти интегралы существуют как абсолютно сходящиеся (а не только в смысле главного значения). Термины «преобразование» и «обратное преобразование», а также символика $F[f]$ и $F^{-1}[f]$ оправдываются следующей теоремой.
\end{notion}

\begin{thm}[об обратимости преобразования Фурье]
\label{th25:3}
Пусть функция $f \in L_R(-\infty, +\infty)$ в каждой точке непрерывна и имеет конечные односторонние производные. Тогда
$$
F^{-1}[F[f]] = f, \quad F[F^{-1}[f]] = f.
$$
\end{thm}


\begin{notion}
Если интеграл Фурье является непрерывным аналогом тригонометрического ряда Фурье, то преобразование Фурье является непрерывным аналогом коэффициентов Фурье в комплексной форме. В прикладных науках последовательность коэффициентов Фурье часто называют \textit{дискретным преобразованием Фурье функции}.
\end{notion}

\section{Преобразование Фурье производной и производная преобразования Фурье}

\begin{thm}[дифференцирование по параметру]
\label{th24:6}
Пусть функция двух переменных $f(x, \alpha)$ непрерывна на множестве
$$
X = \{(x, \alpha) : a \le x < b,\ A \le \alpha \le B\},
$$
где $b$~--- конечно или $+\infty$, вместе со своей частной производной $\dfrac{\partial f}{\partial \alpha}(x, \alpha)$ (имеется в виду, что $\dfrac{\partial f}{\partial \alpha}$ продолжается с $\operatorname{int} X$ на $X$ до функции, непрерывной на $X$).

Пусть далее $\int_a^{\to b} \dfrac{\partial f}{\partial \alpha}(x, \alpha)\,dx$ сходится равномерно по $\alpha \in [A;B]$, а интеграл $I(\alpha) = \int_a^{\to b} f(x, \alpha)\,dx$ сходится при $\alpha = A$.

Тогда интеграл $I(\alpha)$ сходится при всех $\alpha \in [A;B]$; при этом функция $I(\alpha)$ непрерывно дифференцируема на $[A;B]$ и для всех $\alpha \in [A;B]$ выполняется равенство $I'(\alpha) = \int_a^{\to b} \dfrac{\partial f}{\partial \alpha}(x, \alpha)\,dx$ (в точках $A$ и $B$ производная односторонняя), т.е.
$$
\frac{d}{d\alpha}\left(\int_a^{\to b} f(x, \alpha)\,dx\right) = \int_a^{\to b} \frac{\partial f}{\partial \alpha}(x, \alpha)\,dx.
$$
\end{thm}


\begin{thm}[производная преобразования Фурье]
\label{th25:5}
Если функция $f$ непрерывна на $(-\infty; +\infty)$, причём $f(x)$ и $xf(x)$ абсолютно интегрируемы на $(-\infty; +\infty)$, то преобразование Фурье $f$~--- непрерывно дифференцируемая функция на $(-\infty; +\infty)$ и
$$
\frac{d}{dy}(F[f](y)) = F[-ixf(x)].
$$
\end{thm}

\begin{proof}
Дифференцируя по параметру $y$ интеграл
$$
\widehat{f}(y) = \frac{1}{\sqrt{2\pi}}\int_{-\infty}^{+\infty} f(x)e^{-ixy}\,dx,
$$
получим
$$
\frac{d}{dy}\widehat{f}(y) = \frac{1}{\sqrt{2\pi}}\int_{-\infty}^{+\infty} (-ix)f(x)e^{-ixy}\,dx = F[-ixf(x)].
$$

Для обоснования дифференцирования интеграла $\widehat{f}(y)$ по параметру при любом значении $y$ заметим, что сам интеграл $\widehat{f}(y)$ сходится при любом $y$, а $xf(x) \in L_R(-\infty, +\infty)$, поэтому из равенства $|-ixf(x)e^{-ixy}| = |xf(x)|$, выполненного при всех $y$, и признака Вейерштрасса равномерной сходимости несобственных интегралов следует, что продифференцированный интеграл сходится равномерно по $y \in \bbR$. Теорема~\ref{th24:6} применима для $y \in [A;B]$ на любом конечном отрезке $[A;B]$, а значит, нужное равенство имеет место при всех $y \in \bbR$.
\end{proof}

\begin{cons}
Если функция $f$ непрерывна на $(-\infty; +\infty)$, причём при некотором $n = 1, 2, \dots$ функции $f(x), xf(x), \dots, x^n f(x)$ абсолютно интегрируемы на $(-\infty; +\infty)$, то преобразование Фурье $f$ имеет $n$ непрерывных производных на $(-\infty; +\infty)$ и
$$
\frac{d^k}{dy^k}(F[f](y)) = F[(-ix)^k f(x)], \quad k = 1, 2, \dots, n.
$$
\end{cons}

\begin{proof}
Применить $k$ раз Теорему~\ref{th25:5}.
\end{proof}

\begin{defn}
Функцию $f$, определённую на промежутке $I$, будем называть \textit{кусочно-гладкой}\rindex{функция!кусочно-гладкая}, если она кусочно непрерывно дифференцируема на $I$.
\end{defn}

\begin{thm}
\label{th13:9}
Если функция $f$ интегрируема на любом отрезке $[a;b']$, где $b' > a$, и $\lim\limits_{x \to +\infty} f(x) = C > 0$ или $\lim\limits_{x \to +\infty} f(x) = +\infty$, то $\int_{a}^{+\infty} f(x)\,dx$ расходится.
\end{thm}

\begin{thm}[преобразование Фурье производной]
\label{th25:6}
Пусть функция $f$ является кусочно-гладкой\rindex{функция!кусочно-гладкая} на любом конечном отрезке, причём $f \in L_R(-\infty, +\infty)$ и $f' \in L_R(-\infty, +\infty)$. Тогда при всех $y \in \bbR$
$$
F[f'](y) = iy F[f](y).
$$
\end{thm}

\begin{proof}
Так как функция $f$ является кусочно-гладкой на любом конечном отрезке, то $\int_0^x f'(t)\,dt = f(x) - f(0)$ (производная $f'$ кусочно-непрерывна на $[0;x]$ и для неё выполняется формула Ньютона"--~Лейбница~--- Теорема~\ref{cons:newton_leibniz}).

Из сходимости $\int_0^{+\infty} f'(t)\,dt$ следует существование конечного предела $\lim\limits_{x \to +\infty} f(x) = C$. Но если $C \ne 0$, то $\int_0^{+\infty}|f(x)|\,dx$ расходится (Теорема~\ref{th13:9} в применении к функции $|f(x)|$). Значит, $\lim\limits_{x \to +\infty} f(x) = 0$; аналогично $\lim\limits_{x \to -\infty} f(x) = 0$.

Тогда, интегрируя по частям, получим
\begin{multline*}
F[f'] = \frac{1}{\sqrt{2\pi}}\int_{-\infty}^{+\infty} f'(x)e^{-ixy}\,dx = \\
= \frac{1}{\sqrt{2\pi}} f(x)e^{-ixy}\bigg|_{-\infty}^{+\infty}
+ \frac{iy}{\sqrt{2\pi}}\int_{-\infty}^{+\infty} f(x)e^{-ixy}\,dx = iy F[f],
\end{multline*}
так как $|e^{-ixy}| = 1$ и $\lim\limits_{x \to \infty} f(x)e^{-ixy} = 0$.
\end{proof}

\begin{cons}
Пусть функция $f$ такова, что при некотором $n = 1, 2, \dots$ функция $f^{(n-1)}(x)$ является кусочно-гладкой на любом конечном отрезке, причём $f, f', \dots, f^{(n-1)}, f^{(n)} \in L_R(-\infty, +\infty)$. Тогда при всех $y \in \bbR$
$$
F[f^{(n)}](y) = (iy)^n F[f](y),
$$
следовательно,
$$
F[f](y) = o\!\left(\frac{1}{y^n}\right), \quad y \to \infty.
$$
\end{cons}

\begin{proof}
Применить $n$ раз Теорему~\ref{th25:6}. Оценка $o(1/y^n)$ для $F[f](y)$ следует из Теоремы~\ref{th25:4}.
\end{proof}

\begin{notion}[Выводы]
Чем быстрее функция $f$ стремится к нулю при $x \to \infty$ (точнее, чем при большем значении $n$ абсолютно сходится $\int_{-\infty}^{+\infty} x^n f(x)\,dx$), тем более гладкой функцией является преобразование Фурье $F[f](y)$. Наоборот, чем более гладкой является функция $f$, тем быстрее стремится к нулю при $y \to \infty$ её преобразование Фурье.
\end{notion}